\section{Discussion}
\label{sec:discussion}

The theoretical result is a distinction between targets under an explicit observation model. With fixed, known reporting, exogenous exposure and conditionally independent counting errors, current-value risk is $A/m$ and risk for the realised mean of observed periods is $A/(Tm)$. Estimating a stationary process mean additionally incurs the variance of that realised mean. This contribution is unchanged when cases increase within fixed periods, but its behaviour as periods accumulate depends on temporal covariance. Reporting volatility changes even the current and realised-mean risks. These identities provide a reference for interpreting case volume; their assumptions determine where that reference applies.

The variance structure has precedents in demographic and environmental stochasticity and in clustered sampling \citep{engen1998demographic, cochran1977sampling}. Its surveillance interpretation depends on what varies together and what the estimate is intended to describe. Independent offspring and independent reporting generate sampling error that decreases with exposure. Shared transmission conditions affect a process-mean estimate through the selection of periods. Shared reporting fluctuations affect the measurement of those periods. Calling every persistent component an environmental floor obscures this distinction and can mistakenly treat variation in the target as measurement error.

The conditional information calculation complements analyses of information in epidemic curves \citep{parag2022quantifying, parag2022fundamental}. It establishes efficiency of a scaled count in a specified unbiased estimator class with a known denominator and reporting law. It does not establish a universal lower bound for reproduction-number inference. Additional observations, independent reporting measurements and prior structure change the information available. Procedures that trade conditional bias for reduced variance also fall outside that class. A balance scale should therefore always be attached to its estimator, target and sampling experiment rather than interpreted as an intrinsic property of a pathogen.

A related distinction concerns estimation windows. In the Cori framework, widening a window accumulates infectiousness and incidence while imposing local constancy of the reproduction number over a longer interval \citep{cori2013new, gostic2020practical}. This trades temporal resolution against precision and can introduce mismatch when transmission changes. It is not equivalent to drawing more independent transmission periods, nor does it preserve the estimand automatically. The repeated-period averages in Table~\ref{tab:targets} make a different assumption: each period has its own transmission value and a conditionally unbiased measurement. Applying their formulas to overlapping windows requires explicit treatment of overlap and temporal dependence.

The empirical analysis concerns squared relative deviations of apparent weekly statistics, not mean squared errors against observed $R_t$. The affine and power-law comparisons in Section~\ref{sec:res-empirical} address the shape supported over the observed count range. A positive fitted intercept by itself does not establish a nonzero asymptote, and predictive similarity does not identify the mechanism generating either curve. Training-only centring and scoring prevent information leakage in predictive comparisons. Prospective curve-discrimination exercises remain conditional on their generators and address a different question from identifying biological transmission variance.

Controlled thinning adds a known sampling operation to recorded counts. It preserves the underlying observed trajectory before random removal, but nonlinear transformations, continuity corrections and random denominators can change both fitted coefficients. The derivative-based reference is an approximation to this pipeline, not an exact invariant of the statistic. Section~\ref{sec:res-thinning} must therefore be interpreted alongside the retention-specific responses, uncertainty and known-count-model checks. Agreement with a reference response would support that local sampling description; it would not validate the full $A/m$ identity for admissions or identify the origin of an unthinned slope. Failure to reject an intercept change is not evidence of equivalence.

State-intercept fits and alternative centring schemes answer separate questions. State intercepts allow baseline loss to differ across states while estimating an association with within-state count variation. They do not make that variation exogenous or remove common weekly shocks. Resampling summaries must match the intended state distribution, including sampled multiplicities; absent states have no estimated intercept. Recentring by season, trend or national week changes the loss being modelled. The resulting intercept changes are sensitivity measures, not orthogonal biological variance components. Any variance ratio for national-week means concerns signed deviations as defined in Methods, not automatically the squared-loss response.

Spatial and temporal dependence also constrain uncertainty. Regions observed in the same week share influences, while repeated observations within a region share history. Crossed resampling addresses these two sources of dependence under its own assumptions \citep{owen2007pigeonhole, cameron2011robust}; restricted time-block sensitivity explores residual serial dependence. Neither additional rows within the same season nor increasing admissions supplies independent seasons. Similarly, spatial aggregation can change the variability of transmission itself through covariance between local epidemics, as Eq.~(\ref{eq:pooling}) shows. A cross-sectional count association can consequently combine sampling precision with differences in the populations being averaged.

Simulation-based inversion is conditional on the transmission and observation mechanism (Sections~\ref{sec:meth-inversion} and~\ref{sec:res-sim}). The weekly reproduction-number variation specified in a daily renewal simulator is a model parameter indexed by calendar week. It is not directly observed per-generation truth, and multiplying case counts by a time conversion does not convert its variance or a balance scale. Independent observation dispersion, shared reporting shocks and transmission persistence must be specified separately. Agreement between statistics computed from the same admissions is a useful consistency check, but shared data and calibration assumptions limit its independence as evidence.

Generation intervals further condition this mapping. Serial intervals, intrinsic generation intervals and realised household generation intervals measure different quantities \citep{cowling2009estimation, vink2014serial, chan2025estimating}. A sensitivity distribution with an assumed coefficient of variation must be distinguished from the source study's reported spread. Likewise, confidence limits for an interval mean do not define a posterior distribution for the entire generation interval. Conditional calibration intervals and changes across specified scenarios describe different uncertainties; a scenario range is not a joint confidence interval. Propagating uncertainty requires a justified distribution and an explicit joint model, not a relabelling of sensitivity endpoints.

Finally, facility participation measures only one aspect of ascertainment. Testing, care seeking, coding and admission practice can vary within participating hospitals. Comparisons across reporting regimes remain descriptive without a design identifying a policy effect. Independent measurements of these processes, finer spatial observations and further seasons would directly address the principal identification limits. The symbolic design extensions in Appendix~\ref{app:extensions} illustrate how a chosen target enters allocation, subject to costs, feasible periods and dependence. Operational choices require those quantities to be measured on consistent scales. The empirical summaries alone do not supply a universal case threshold or a uniquely identified transmission variance.
